\documentclass[12pt,a4paper]{article}
\usepackage[margin=1in]{geometry}
\usepackage[T1]{fontenc}
\usepackage[utf8]{inputenc}
\usepackage{graphicx}
\usepackage{xcolor}
\usepackage{booktabs}
\usepackage{hyperref}
\usepackage{float}
\usepackage{caption}
\usepackage{tocloft}
\usepackage{listings}

\setlength{\cftbeforefigskip}{0pt}
\lstset{
  basicstyle=\ttfamily\small,
  language=Python,
  keywordstyle=\color{blue},
  commentstyle=\color{gray},
  stringstyle=\color{teal},
  showstringspaces=false,
  breaklines=true,
  frame=single,
  framerule=0.4pt,
  rulecolor=\color{gray}
}

\title{Relay Network Simulator\\Documentatie Tehnica}
\author{}
\date{}

\begin{document}

\pagenumbering{gobble}
\begin{titlepage}
  \centering
  \vspace*{1.5cm}
  {\LARGE \textbf{Relay Network Simulator}\par}
  \vspace{0.6em}
  {\Large Documentatie Tehnica\par}
  \vspace{0.8em}
  {\color{gray}\rule{0.7\textwidth}{0.6pt}}\par
  \vspace{2em}
  \begin{tabular}{rl}
    \textbf{Autor:} & Jalbă Simion \\
    \textbf{Materie:} & Rețele de calculatoare \\
    \textbf{Universitate:} & UTCN \\
    \textbf{Facultate:} & Automatică si Calculatoare \\
  \end{tabular}
  \vfill
  {\color{gray}\rule{0.5\textwidth}{0.4pt}}\par
  \vspace*{0.8cm}
\end{titlepage}

\clearpage
\pagenumbering{roman}
\tableofcontents
{\small
\listoffigures
}
\clearpage
\pagenumbering{arabic}

\section{Introducere}
Acest document descrie proiectul Relay Network Simulator, un simulator de propagare a mesajelor in retele distribuite. Proiectul urmareste compararea a doua protocoale de diseminare (Flooding si Gossip) in diferite topologii, cu si fara pierderi de pachete.

\section{Obiective}
\begin{itemize}
  \item Modelarea retelei ca graf si simularea propagarii mesajelor.
  \item Evaluarea compromisului dintre acoperire si overhead.
  \item Generarea de rezultate reproductibile prin seed-uri controlate.
\end{itemize}

\section{Structura proiectului}
Proiectul este organizat modular:
\begin{itemize}
  \item \textbf{config}: parametri globali si configurari.
  \item \textbf{experiments}: scenarii si rularea experimentelor.
  \item \textbf{metrics}: colectare si agregare rezultate.
  \item \textbf{network}: modelul de graf si elementele de baza.
  \item \textbf{protocols}: implementari Flooding si Gossip.
  \item \textbf{simulation}: motorul de simulare si scheduler.
  \item \textbf{visualization}: grafice si vizualizari.
\end{itemize}

\section{Arhitectura si flux de date}
Simularea este orchestrata din \texttt{main.py}. Fiecare experiment construieste o topologie, ruleaza simularea pentru un protocol si colecteaza metrici. Rezultatele sunt salvate in CSV si vizualizari PNG.

\begin{figure}[H]
  \centering
  \includegraphics[width=0.85\textwidth]{C:/Users/jalba/Desktop/RC/miniproiect/relay-network-simulator/results/plots/experiment1_topology.png}
  \caption{Exemplu de topologie generata pentru simulare.}
\end{figure}

\section{Modelul de simulare}
Reteaua este reprezentata ca graf (NetworkX). Simulatorul foloseste un scheduler de evenimente discrete pentru a simula transmiterea mesajelor intre noduri. Se pot introduce pierderi de pachete prin parametrul \texttt{loss\_rate}.

\begin{lstlisting}
python main.py --experiment all --seed 42
\end{lstlisting}

\begin{lstlisting}
sim_cfg = SimulationConfig(loss_rate=loss_rate, seed=seed)
sim = RelayNetworkSimulator(graph=graph, protocol=protocol, config=sim_cfg)
result = sim.run(source_id=0)
\end{lstlisting}

\section{Protocoale}
\subsection{Flooding}
Fiecare nod retransmite mesajul catre toti vecinii sai. Asigura acoperire ridicata, dar produce overhead mare.

\subsection{Gossip}
Fiecare nod retransmite mesajul catre un subset aleator de vecini (parametrul \texttt{k}). Reduce overhead-ul, dar poate reduce acoperirea.

\begin{lstlisting}
protocol = FloodingProtocol() if protocol_name == "flooding" else GossipProtocol(k=2)
\end{lstlisting}

\begin{lstlisting}
graph = TopologyBuilder.create(
  topology,
  node_count=node_count,
  seed=seed,
  **topology_kwargs,
)
\end{lstlisting}

\section{Experimente}
\subsection{Experiment 1: Comparatie protocol fara packet loss}
Comparatie directa intre Flooding si Gossip in topologii random si small-world.

\subsection{Experiment 2: Sweep de pierderi}
Se variaza \texttt{loss\_rate} si se observa impactul asupra delivery ratio.

\subsection{Experiment 3: Scalare dupa numar de noduri}
Se analizeaza performanta pentru retele de 50, 100 si 500 de noduri.

\subsection{Experiment 4: Topologii rare vs dense}
Se compara protocoalele pentru grafuri random cu densitate mica si mare.

\begin{lstlisting}
for protocol_name in ("flooding", "gossip"):
    for density_label, edge_probability in density_map.items():
        record = _run_single(
            protocol_name=protocol_name,
            topology="random",
            node_count=cfg.default_node_count,
            loss_rate=0.1,
            seed=run_seed,
            topology_kwargs={"edge_probability": edge_probability},
        )
\end{lstlisting}

\begin{lstlisting}
ChartBuilder.packet_loss_curve(
    df,
    "delivery_ratio",
    cfg.results_plots_dir / "experiment2_packet_loss_delivery.png",
    "Experiment 2: Delivery Ratio",
)
\end{lstlisting}

\section{Metrici}
Metricile principale colectate sunt:
\begin{itemize}
  \item \texttt{delivery\_ratio}
  \item \texttt{network\_overhead}
  \item \texttt{total\_transmissions}
  \item \texttt{avg\_hop}, \texttt{avg\_delay}
  \item \texttt{loss\_impact}
\end{itemize}

\begin{lstlisting}
return {
  "delivery_ratio": delivery_ratio,
  "total_transmissions": float(self.total_transmissions),
  "redundant_messages": float(self.redundant_messages),
  "avg_hop": avg_hop,
  "avg_delay": avg_delay,
  "network_overhead": overhead,
  "loss_impact": loss_impact,
  "unique_receivers": float(unique_receivers),
}
\end{lstlisting}

\section{Rezultate si grafice}
\begin{figure}[H]
  \centering
  \includegraphics[width=0.85\textwidth]{C:/Users/jalba/Desktop/RC/miniproiect/relay-network-simulator/results/plots/experiment1_delivery_ratio.png}
  \caption{Delivery ratio pentru Flooding vs Gossip (Experiment 1).}
\end{figure}

\begin{figure}[H]
  \centering
  \includegraphics[width=0.85\textwidth]{C:/Users/jalba/Desktop/RC/miniproiect/relay-network-simulator/results/plots/experiment1_overhead.png}
  \caption{Network overhead pentru Flooding vs Gossip (Experiment 1).}
\end{figure}

\begin{figure}[H]
  \centering
  \includegraphics[width=0.85\textwidth]{C:/Users/jalba/Desktop/RC/miniproiect/relay-network-simulator/results/plots/experiment2_packet_loss_delivery.png}
  \caption{Impactul pierderilor asupra delivery ratio (Experiment 2).}
\end{figure}

\begin{figure}[H]
  \centering
  \includegraphics[width=0.85\textwidth]{C:/Users/jalba/Desktop/RC/miniproiect/relay-network-simulator/results/plots/experiment3_transmissions.png}
  \caption{Total transmissions pentru dimensiuni de retea diferite (Experiment 3).}
\end{figure}

\begin{figure}[H]
  \centering
  \includegraphics[width=0.85\textwidth]{C:/Users/jalba/Desktop/RC/miniproiect/relay-network-simulator/results/plots/experiment4_overhead.png}
  \caption{Overhead in topologii rare vs dense (Experiment 4).}
\end{figure}

\section{Concluzii}
Simulatorul evidentiaza compromisurile intre acoperire si costul de comunicare. Flooding obtine, in general, delivery ratio mai mare, dar cu overhead ridicat. Gossip reduce overhead-ul, dar poate pierde acoperire, mai ales la pierderi mari.

\end{document}
